% TOP_PROBLEM: {"id": 3121, "title": "Minimal graphs with a prescribed number of spanning trees", "category_id": 3, "category": {"id": 3, "name": "graph_theory", "display_name": "Graph Theory", "description": "Problems involving graphs, networks, and their properties.", "slug": "graph-theory", "order_index": 3, "created_at": "2026-07-31T15:26:25.671Z"}, "status": "open", "problem_number": "OPG-37420", "set_id": 12}
% FIRST_POSTED: 2026-10-01
% ATTEMPT_STATUS: unresolved
% UnsolvedMath ID 3121; source catalogue code OPG-37420
% !TeX program = lualatex
% GPT-6 Astra Ultra research attempt, 2026-10-01; self-contained partial work.
% Completion estimate: no numerical percentage assigned; full problem unresolved.
\documentclass[11pt,a4paper]{article}
\usepackage[margin=25mm]{geometry}
\usepackage{fontspec}
\setmainfont{lmroman10-regular.otf}[BoldFont=lmroman10-bold.otf,
 ItalicFont=lmroman10-italic.otf,BoldItalicFont=lmroman10-bolditalic.otf]
\usepackage{amsmath,amssymb,mathtools}
\usepackage{unicode-math}
\setmathfont{latinmodern-math.otf}
\AtBeginDocument{\let\setminus\smallsetminus}
\usepackage{xurl,hyperref}
\hypersetup{colorlinks=true,urlcolor=blue}
\setcounter{secnumdepth}{0}
\setlength{\parindent}{0pt}
\setlength{\parskip}{5pt}
\setlength{\emergencystretch}{3em}
\begin{document}
\section*{Minimal graphs with a prescribed number of spanning trees}\label{3121}
{\small UnsolvedMath ID 3121 · OPG-37420 · Graph Theory · OpenGarden}
\subsection{Definitions and mathematical statement}
For a finite simple undirected graph $G$, a spanning tree is a connected
acyclic subgraph containing every vertex. Let $\tau(G)$ be the number
of its spanning trees, with distinct edge sets counted separately.
For each integer $n\ge3$, define
\[
 \alpha(n)=\min\{|V(G)|:G\text{ finite, simple and undirected},\
                                  \tau(G)=n\}.
\]
The minimum exists: deleting any one edge of the cycle $C_n$ gives
one of its exactly $n$ spanning trees, so $\alpha(n)\le n$.
Any graph contributing to the minimum is connected because its tree
count is positive. Loops, parallel edges and directed arcs are not allowed.

The conjecture asserts
\[
 \lim_{n\to\infty}\frac{\alpha(n)}{\log n}=0,
\]
where $n$ tends to infinity through all integers and $\log$ may be taken
to be the natural logarithm. Equivalently, for every $\varepsilon>0$
there exists $N_\varepsilon$ such that every integer $n\ge N_\varepsilon$
is the exact spanning-tree count of a simple graph on at most
$\varepsilon\log n$ vertices.

The number $n$ here is the prescribed count, not the number of graph
vertices. The conclusion is uniform over all sufficiently large counts;
constructions only for a sparse subsequence of highly composite integers
do not suffice. Nor is it enough to find a graph with at least $n$
spanning trees. The exact equality is what makes the minimization an
arithmetic realization problem.
\subsection{Short English statement}
Can every sufficiently large integer be realized as an exact spanning-tree count using a number of vertices negligible compared with its logarithm?
\subsection{Research attempt}
\paragraph{Scope and process.}
GPT-6 Astra Ultra, 1 October 2026. This attempt keeps the catalogue statement above, checks primary literature, and develops a restricted argument with an adversarial gap audit. The proofs below are the mathematical output of this attempt, not a claim of novelty or external certification.

\paragraph{Literature and attack plan.}
Stong's primary paper [R1] proves the uniform upper bound $O((\log n)^{3/2}/\log\log n)$ and explicitly distinguishes it from the conjectured $o(\log n)$. We do not claim that even an $O(\log n)$ bound follows from that result. The attempt here determines an exact infinite subsequence, proves the universal information-size obstruction, and examines why gluing graphs fails to supply the required uniform little-oh bound.

\paragraph{A universal counting obstruction.}
A simple graph on $v\ge2$ labelled vertices has at most $v^{v-2}$ spanning trees, since each of its trees is also a tree of $K_v$, and Prüfer coding counts the latter. For completeness, the encoding successively removes the smallest labelled leaf and records its neighbour until two vertices remain. It produces a word of length $v-2$ over the $v$ labels. Decoding chooses at each step the smallest label absent from the remaining word, joins it to the first word entry, and removes that entry; finally it joins the last two labels. These procedures are inverse, so exactly $v^{v-2}$ labelled trees exist.

Thus every graph realizing $n\ge3$ satisfies
\[
 \log n\le(v-2)\log v. \tag{1}
\]
For sufficiently large $n$, write $L=\log n$. If $v\ge L$, then $v\ge L/\log L$. If $v<L$, (1) gives $L\le v\log L$, giving the same inequality. Consequently
\[
 \alpha(n)\ge\frac{\log n}{\log\log n}
 \quad\text{for all sufficiently large }n. \tag{2}
\]
This lower bound is compatible with $o(\log n)$; its ratio to $\log n$ tends to zero.

\paragraph{An exact subsequence.}
For $v\ge3$, put $n_v=v^{v-2}$. The complete graph gives $\alpha(n_v)\le v$. No graph of smaller order can attain that count: for integers $2\le u<v$, $u^{u-2}<v^{v-2}$, as follows by increasing both positive base and nonnegative exponent, with strictness once $v\ge3$. Graphs on one or two vertices have at most one spanning tree. Therefore
\[
 \alpha(v^{v-2})=v. \tag{3}
\]
It follows both that
\[
 \frac{\alpha(n_v)}{\log n_v}
 =\frac{v}{(v-2)\log v}\longrightarrow0
\]
and that the lower-bound scale in (2) is attained asymptotically:
\[
 \liminf_{n\to\infty}
 \frac{\alpha(n)\log\log n}{\log n}=1. \tag{4}
\]
For the upper direction of (4), use (3) and
$\log\log(v^{v-2})=\log(v-2)+\log\log v$; division by $\log v$ tends to one. The lower direction is (2).

\paragraph{Exact multiplication by vertex gluing.}
Let connected simple graphs $G,H$ share exactly one vertex and otherwise have disjoint vertices and edges. A spanning tree of their union restricts to a spanning tree on each graph. Indeed, a path from a vertex of $G$ to the shared vertex cannot leave $G$ before first reaching that vertex; connectedness of the whole tree therefore forces connectedness of the restriction, while acyclicity is inherited. Conversely the union of two spanning trees is connected and acyclic, since a cycle crossing between the pieces would have to revisit the unique shared vertex. Thus
\[
 \tau(G\vee H)=\tau(G)\tau(H),\qquad
 \alpha(ab)\le\alpha(a)+\alpha(b)-1 \quad(a,b\ge3). \tag{5}
\]
Iterating gives $\alpha(a^t)\le1+t(\alpha(a)-1)$.

\paragraph{Why the gluing approach stops.}
For a fixed factor $a$, division of the last inequality by $\log(a^t)=t\log a$ only bounds the ratio by a positive constant. Letting the efficient block size grow can help selected factorizations, but an arbitrary integer need not admit such factors; primes admit none at all. Neither the complete-graph subsequence nor (5) turns into a construction for every sufficiently large integer. Changing ``exactly $n$'' to ``at least $n$'' would remove the arithmetic difficulty and change the problem.

\paragraph{Verification and final status.}
The constructions remain simple graphs: identifying one vertex introduces neither loops nor parallel edges. For $v=3$, (3) gives the triangle with three trees; for $v=4$, it gives $K_4$ with sixteen trees. The limiting calculation in (4) concerns a liminf, not a limit or a uniform upper bound. The exact subsequence, product identity, and obstruction (2) are proved. The all-integer realization bound in the original conjecture is \textbf{UNRESOLVED} by this attempt.

\subsection{Sources}
\begin{itemize}
\item[S1] ULAM.AI and the UnsolvedMath Contributors, supplied problems.json, record 3121 (OPG-37420), OpenGarden collection. Catalogue identity and source transcription; CC BY 4.0.
\url{https://huggingface.co/datasets/ulamai/UnsolvedMath}.
\item[S2] Open Problem Garden, Minimal graphs with a prescribed number of spanning trees. Original problem and discussion; the mathematical formulation below is expanded from the supplied source record.
\url{http://www.openproblemgarden.org/op/minimal_graphs_with_a_prescribed_number_of_spanning_trees}.
\item[R1] R. Stong, \emph{Minimal graphs with a prescribed number of spanning trees}, Australasian Journal of Combinatorics 82(2) (2022), 182--196. Primary abstract and introduction checked.
\url{https://ajc.maths.uq.edu.au/pdf/82/ajc_v82_p182.pdf}.
\end{itemize}
\end{document}
